\subsection{\texorpdfstring{THE TOPOLOGY OF \(\mathbf{R}^n\)}{THE TOPOLOGY OF R TO THE N}}

DEFINITION 1. The topological product of n spaces identical with the real line is called real number space of n dimensions, and is denoted by \(R^{n}\) .

Remark. The space \(R^{0}\) consists of a single point.

From Set Theory, Chapter III, § 6, no. 3, Theorem 2, Corollary 1, we know that, if E is an infinite set, \(E^{n}\) is equipotent with E for all integers n \textgreater{} 0; hence, if n \textgreater{} 0, \(R^{n}\) is equipotent with R, that is, \(R^{n}\) has the power of the continuum (cf.~Exercises 1 and 2).

DEFINITION 2. Any subset of \(\mathbf{R}^n\) which is the product of \(n\) open (resp. closed) intervals of \(\mathbf{R}\) is called an open (resp. closed) box in \(\mathbf{R}^n\) . {[}For \(n = 2\) it is called an open (resp. closed) rectangle.{]}

The open boxes in \(\mathbf{R}^n\) form a base of the topology of \(\mathbf{R}^n\) (Chapter I, § 4, no. 1); the open boxes which contain a point \(x = (x_i)_{1 \leqslant i \leqslant n}\) of \(\mathbf{R}^n\) form a fundamental system of neighbourhoods of \(x\) , and so do the closed boxes of \(\mathbf{R}^n\) for which \(x\) is an interior point.

Every non-empty open box in \(\mathbf{R}^n\) is homeomorphic to \(\mathbf{R}^n\) (Chapter IV, § 4, no. 1, Proposition 1).

It follows that, when \(n \geqslant 1\) , every non-empty open set in \(R^{n}\) has the power of the continuum.

An open (resp. closed) cube of \(R^{n}\) is an open (resp. closed) box which is the product of n bounded intervals of equal length {[}for n = 2, it is called an open (resp. closed) square{]}; the common length of these intervals is called the side (or side-length) of the cube. The open cubes

\[
\mathrm{K} _ {m} = \prod_ {1 \leqslant i \leqslant n} \left] x _ {i} - \frac {1}{m}, x _ {i} + \frac {1}{m} \right[
\]

form a countable fundamental system of neighbourhoods of the point \(\boldsymbol{x} = (x_{i})\) , as m runs through the set of all integers \textgreater0 or through any sequence of integers increasing to infinity.

Every open (or closed) box in \(\mathbf{R}^n\) is connected (Chapter I, § 11, no. 4, Proposition 8); in particular, \(\mathbf{R}^n\) is connected and locally connected.

If A is a non-empty open set in \(R^{n}\) , its components are therefore open sets (Chapter I, §11, no. 6, Proposition 11); and the set of these components is countable, for \(R^{n}\) has a countable dense subset (for example \(Q^{n}\) ).

Consider under what conditions a subset A of \(R^{n}\) will be relatively compact. By Tychonoff's theorem (Chapter I, § 9, no. 5, Theorem 3) it is necessary and sufficient that the projections of A on the factors of \(R^{n}\) should be relatively compact; by the Borel-Lebesgue theorem (Chapter IV, § 2, no. 2, Theorem 2) this is equivalent to saying that these projections are bounded subsets of R. We say that a subset A of \(R^{n}\) is bounded if all its projections are bounded subsets of R; thus we have proved:

PROPOSITION 1. A subset A of \(R^{n}\) is relatively compact if and only if it is bounded.

COROLLARY. The space \(\mathbf{R}^n\) is locally compact, but is not compact if \(n \geqslant 1\) .

